\section{病理的描像と限界}
\label{sec:pathologies}

\subsection{意図的に除外するもの}
本枠組みは以下を意図的に除外する。

\textbf{意味理論。}
指示、真理条件、解釈内容については主張しない。
意味という語が現れる場合、それは安定した軌道に適用される創発的な呼称であり、理論の原始概念ではない。

\textbf{現象学と意識。}
本枠組みは主観的経験を生成せず、説明もしない。
意識的アクセスは、現在repository内にsourceが存在するConsciousness Theoryとの比較に向けた、下流の主題として扱う。
この比較は、CTに機構を移し入れるものでも、その現象的解釈を検証するものでもない。

\textbf{神経生物学的実装。}
本理論は、特定の神経アーキテクチャに依存しない抽象度で、構造的・動的条件を記述する。
神経実装との対応は今後の課題とする。

\subsection{失敗領域の提案}
以下は定性的な領域仮説であり、診断的・臨床的分類器でも、認知状態の必要条件や十分条件でもない。

\textbf{過閉包（深すぎる井戸）。}
$\sigma_{\mathrm{drive}}$に比べて有効な$\sigma_{\mathrm{barrier}}$が大きいことを、地形変化に抵抗する領域として解釈する。
提案する認知的対応は、より反復的で適応性の低い思考である。
硬直性、教条性、定着した信念は、ここでは現象学的な比較の語であり、パラメータによって同定または因果的に説明された状態ではない。

\textbf{閉包不足（浅すぎる井戸）。}
$\sigma_{\mathrm{dissip}}$の優勢を、モデル描像では井戸の形成より速い浅化と関連づける。
これに不安定な思考や概念の安定化の困難さとの対応を提案する。
理解の持続性の低さとの関係は定性的な仮説であり、持続を阻む機構として確立されたものではない。

\textbf{同期の破綻（社会的スケール）。}
事前の地形の乖離、共有された接触の不足、意図的な操作は、エージェント間の同期を弱めると提案する。
共有地形の断片化は、このパターンのモデル解釈であり、コミュニケーションの失敗、文化的対立、認識的分極化との比較候補を与える。
これらの結果を一つの測定された領域と同一視したり、集合的幾何だけから推論したりするものではない。

\subsection{対象範囲の限界}
本理論は構造的・動的記述を提案するが、以下を確立しない。
\begin{itemize}
\item 特定の認知系についての定量的予測（そのためには$K$、$D$、$\mu$、$\alpha$、$\Gamma$、$\kappa_B$の具体的指定が必要である）。
\item $\Phi$の幾何と、観測可能な神経・行動データとの対応条件。
\item 言語獲得、意味変化、集合的知性の完全な理論。
\end{itemize}
これらは提案された枠組み内の今後の課題として残る。
